%% kapa: exam=1 round=35 year=2024 period=1 form=A subject=economics subject_ko=경제학원론 seq=5 number=45
%% source: Q-Net 감정평가사 기출문제 (cst00309) artlSeq=5213229 file=raw/1008/5213229/1교시 문제지 원본.pdf
%% conversion: kordoc 4.24.51 -> tools/gen1.py (auto); review=proofread-sonnet; flags=none
\begin{question}{45}
등량곡선에 관한 설명으로 옳은 것을 모두 고른 것은? (단, 한계기술대체율은 절댓값으로 나타낸다.)
\begin{framedtext}
ㄱ. 한계기술대체율은 두 생산요소의 한계생산 비율과 같다. \par
ㄴ. 두 생산요소 사이에 완전 대체가 가능하다면 등량곡선은 직선이다. \par
ㄷ. 등량곡선이 원점에 대해 볼록한 모양이면 한계기술대체율체감의 법칙이 성립한다. \par
ㄹ. 콥-더글러스(Cobb-Douglas) 생산함수의 한계기술대체율은 0 이다. \par
\end{framedtext}
\begin{choices}
  \item ㄱ, ㄴ
  \item ㄴ, ㄷ
  \item ㄷ, ㄹ
  \item ㄱ, ㄴ, ㄷ
  \item ㄱ, ㄴ, ㄷ, ㄹ
\end{choices}
\end{question}
\answer{4}
